\subsection{Spielererfahrung}

\subsubsection{Spielstile}
Die Challenges unterstützen bewusst zwei unterschiedliche Herangehensweisen. 
Spieler können ein Level zunächst analysieren, mögliche Aktionsfolgen gedanklich durchspielen und eine Lösung möglichst weit im Voraus planen. 
Alternativ können sie auch explorativ vorgehen, unmittelbar verschiedene Aktionen ausprobieren und aus deren Auswirkungen lernen. 
Da Challenges beliebig häufig neu gestartet werden können, sind beide Spielweisen gleichermaßen valide und können je nach Präferenz oder Situation miteinander kombiniert werden.
Beide mögliche Spielweisen bleiben auch im Roguelite-Modus erhalten. 
Auch dort kann ein Encounter sorgfältig vorausgeplant werden. 
Es ist aber genauso möglich ihn eher intuitiv aus dem Bauch heraus zu spielen. 
Nicht optimales Spielen führt dabei nicht zu einem kostenlosen Neustart, sondern kann zur Folge haben, dass eine Rast-Aktion erforderlich ist. 
Beide Spielweisen sind in beiden Spielmodi möglich, allerdings belohnt das Roguelite explizit gutes Planen.

\subsubsection{Puzzle Solving als zentrale Spielmotivation}
Problemlösen gehört zu den grundlegenden Motivationen des Spielens. Jesse Schell geht in seiner Definition des Spiels sogar so weit, ein Spiel als „problem-solving activity, approached with a playful attitude“ zu beschreiben. Die Freude daran, Zusammenhänge zu verstehen, eine zunächst schwierige Situation zu durchdringen und schließlich eine Lösung zu finden, bildet entsprechend einen wesentlichen Teil der Spielerfahrung von Cycle of Cinder.
Gleichzeitig reagieren Spieler sehr unterschiedlich auf schwierige Probleme und insbesondere auf wiederholtes Scheitern. Modelle der Goal Orientation unterscheiden beispielsweise im Hinblick auf die Motivation durch Kompetenzerleben zwischen Approach und Avoidance. Vereinfacht bedeutet dies, dass manche Spieler gezielt Situationen suchen, in denen ihre eigene Kompetenz herausgefordert und weiterentwickelt wird, während andere eher vermeiden möchten, mit Situationen konfrontiert zu werden, in denen die Grenzen der eigenen Kompetenz deutlich werden.
Puzzle geben in gewisser Weise ein besonders klares und unmittelbares Feedback über den eigenen Kompetenzstand. Eine Lösung funktioniert oder sie funktioniert nicht. Gerade bei anspruchsvollen Puzzles wird dadurch sehr deutlich erfahrbar, ob die eigene aktuelle Kompetenz ausreicht, um die gestellte Problemstellung zu durchdringen und zu lösen.
Für Spieler mit einer stärkeren Approach-Orientierung kann genau dieses klare Feedback besonders motivierend sein: Die Herausforderung macht sichtbar, dass die eigene Kompetenz noch nicht ausreicht, und erzeugt damit einen konkreten Anlass, sie zu erweitern. Für stärker Avoidance-orientierte Spieler kann dieselbe Klarheit dagegen belastend wirken, weil ein ungelöstes Puzzle sehr unmittelbar die Grenze der eigenen Kompetenz aufzeigt.


Für manche Spieler entsteht ein besonders befriedigendes Erlebnis gerade daraus, zunächst an einer schwierigen Aufgabe zu scheitern, ihre Mechanismen immer besser zu verstehen und sich die Lösung schließlich selbst zu erarbeiten. Für andere kann derselbe langwierige Lösungsprozess dagegen schnell in Frustration umschlagen. Das Design von Cycle of Cinder soll deshalb anspruchsvolles Puzzle Solving ermöglichen, ohne den Fortschritt jedes Spielers zwingend davon abhängig zu machen.

\subsubsection{Challenge-Modus: Herausforderung ohne Sackgasse}
Im Challenge-Modus dürfen einzelne Encounter bewusst anspruchsvoll sein und mehrere Versuche erfordern. Wer die Herausforderung sucht, kann so lange experimentieren und planen, bis die eigene Lösung gefunden ist. Gleichzeitig kann eine Challenge jederzeit übersprungen und später erneut gespielt oder auf Wunsch eine Lösung angezeigt werden.
Damit wird vermieden, dass aus einer interessanten Herausforderung eine dauerhafte Fortschrittsbarriere entsteht. Auch das Anzeigen einer Lösung ist dabei nicht lediglich als „Aufgeben“ zu verstehen: Untersuchungen zu Insight-Problemen zeigen, dass eine präsentierte Lösung insbesondere nach einem zuvor erfolglosen Lösungsversuch selbst ein Aha-Erlebnis und das damit verbundene plötzliche Verständnis auslösen kann.
Roguelite-Modus: Aus Scheitern wird Effizienz
Der Roguelite-Modus löst dasselbe Problem auf eine andere Weise. Hier muss ein Spieler einen Encounter nicht perfekt lösen, um weiterspielen zu können. Wird keine ausreichend effiziente Lösung gefunden, können über die Rest-Action zusätzliche Ressourcen eingesetzt werden. Der Spieler kommt weiter, bezahlt eine weniger gute Lösung jedoch mit einer begrenzten Ressource des Runs – insbesondere mit Zeit.
Dadurch wird aus der binären Frage „Kann ich dieses Puzzle lösen oder nicht?“ eine graduelle Frage: „Wie effizient kann ich diese Situation mit meinem aktuellen Build lösen?“ Ein sehr analytischer Spieler kann weiterhin versuchen, Encounters möglichst optimal vorauszuplanen. Ein stärker intuitiver oder explorativer Spieler kann dagegen ausprobieren, reagieren und bei Bedarf zusätzliche Ressourcen einsetzen.
Auf diese Weise soll der puzzle-artige Gameplay Loop unterschiedliche Präferenzen zulassen, ohne seinen strategischen Anspruch zu verlieren. Spieler, die schwierige Probleme vollständig durchdringen und selbstständig knacken möchten, finden diese Erfahrung insbesondere in den handcrafted Challenges. Im Roguelite-Modus wird dieselbe Fähigkeit zum Problemlösen in einen weniger binären Kontext eingebettet, in dem nicht eine einzelne ungelöste Aufgabe über Erfolg oder Misserfolg entscheidet, sondern die Qualität vieler Entscheidungen über den Verlauf eines gesamten Runs.

\subsubsection{Puzzle geben starkes Feedback}
Puzzle geben ein besonders klares und unmittelbares Feedback über den eigenen Kompetenzstand: Eine Lösung funktioniert oder sie funktioniert nicht. Gerade bei anspruchsvollen Puzzles wird dadurch sehr deutlich erfahrbar, ob die eigene aktuelle Kompetenz ausreicht, um die gestellte Problemstellung zu durchdringen und zu lösen.
Modelle der Goal Orientation unterscheiden im Hinblick auf die Motivation durch Kompetenzerleben unter anderem zwischen Approach und Avoidance. Abstrakt gesprochen suchen manche Spieler gezielt nach Herausforderungen, die ihre Kompetenz ausreizen und ihnen ermöglichen, diese weiterzuentwickeln, während andere eher vermeiden möchten, mit dem Punkt konfrontiert zu werden, an dem ihre aktuelle Kompetenz nicht ausreicht.
Dabei ist für stärker Avoidance-orientierte Spieler nicht zwangsläufig schon ein einzelner Fehler oder ein gescheiterter Lösungsversuch problematisch. Kritischer ist die Situation, über längere Zeit in einem Zustand des Nicht-Lösens festzustecken und keinen erkennbaren Weg zu haben, diesen zu überwinden. Genau hier setzt der Challenge-Modus von Cycle of Cinder an: Encounters können beliebig oft erneut versucht werden, aber auch übersprungen oder ihre Lösung kann auf Wunsch angezeigt werden. Dadurch muss kein Spieler dauerhaft an einer einzelnen Aufgabe festhängen.
Die Möglichkeit unbegrenzter neuer Versuche erfüllt zugleich eine andere Funktion für stärker Approach-orientierte Spieler. Wer eine schwierige Aufgabe gerade deshalb reizvoll findet, weil sie die eigene Kompetenz herausfordert, kann so lange neue Ansätze ausprobieren, bis die Lösung tatsächlich selbst gefunden wurde. Das Spiel unterstützt damit sowohl das Bedürfnis, eine schwierige Herausforderung durch eigene Kompetenz zu überwinden, als auch das Bedürfnis, eine als blockierend empfundene Situation verlassen zu können.

\subsubsection{Roguelite-Struktur als Ausgleich unterschiedlicher Spielstärken}
\todo{
    Andere Option besser?
}
Roguelite-Struktur als Ausgleich unterschiedlicher Spielstärken
Wie schwierig sich ein Spiel für einen Spieler anfühlt, hängt nicht allein von der objektiven Schwierigkeit einer Aufgabe ab. Entscheidend ist vielmehr das Verhältnis zwischen der eigenen Kompetenz und der gestellten Herausforderung. Hinzu kommen weitere Faktoren wie das Feedback des Spiels, die eigenen Erwartungen und der Vergleich mit der bisherigen eigenen Leistung. Ein Spieler kann dieselbe Herausforderung deshalb sehr unterschiedlich wahrnehmen, je nachdem, ob er erwartet hatte, sie problemlos zu bewältigen, ob er bereits zuvor weiter gekommen ist oder ob das Spiel ihm sehr deutlich signalisiert, dass seine aktuelle Kompetenz nicht ausreicht.
Die Roguelite-Struktur von Cycle of Cinder wirkt mehreren dieser Probleme gleichzeitig entgegen. Ein wichtiger Faktor ist, dass der Spieler zu Beginn eines Runs nicht unmittelbar weiß, wie weit er vom endgültigen Ziel entfernt ist oder welchen konkreten Leistungsstand er erreichen müsste, um den gesamten Run erfolgreich abzuschließen. Fortschritt wird zunächst relativ erlebt: Der Spieler erreicht einen neuen Ort, überwindet weitere Encounter, verbessert seinen Build und kommt innerhalb der Reise weiter. Dadurch entsteht weniger häufig die direkte Wahrnehmung, ein zuvor gesetztes Leistungsziel klar verfehlt zu haben.
Gleichzeitig sorgt die persistente Metaprogression dafür, dass sich die Ausgangsbedingungen über mehrere Runs hinweg verbessern und neue strategische Möglichkeiten freigeschaltet werden. Fortschritt entsteht damit nicht ausschließlich dadurch, dass der Spieler selbst schneller lernt oder seine spielerische Kompetenz stark erhöht. Auch Spieler, deren individuelle Lernkurve flacher verläuft, erhalten zunehmend bessere Voraussetzungen und können dadurch im Laufe der Zeit weiter vordringen. Subjektiv entsteht so kontinuierlich das Gefühl von Entwicklung und wachsender Kompetenz, auch wenn dieser Fortschritt teilweise durch das Spielsystem unterstützt wird.
Darüber hinaus reguliert sich die erlebte Schwierigkeit innerhalb eines Runs zu einem gewissen Grad selbst. Solange die Encounter gut zur aktuellen Spielerkompetenz und zum aufgebauten Loadout passen, können sie effizient gelöst werden und der Spieler schreitet schnell voran. Werden die Herausforderungen zunehmend anspruchsvoller, steigt der Ressourcenverbrauch: Es müssen häufiger zusätzliche Ressourcen eingesetzt oder Rest-Actions genutzt werden, wodurch Zeit verloren geht und der Run früher endet. Sehr starke Spieler verbringen entsprechend weniger Zeit mit für sie trivialen Encountern und erreichen schneller schwierigere Bereiche, während weniger erfahrene Spieler länger in einem Schwierigkeitsbereich verbleiben, der ihrer aktuellen Kompetenz entspricht.
Die Roguelite-Struktur erzeugt damit eine Form indirektee dynamische Schwierigkeitsanpassung.

Ein weiterer Vorteil der Roguelite-Struktur ergibt sich aus der inhärenten Varianz des Loadout-Buildings. Welche Zauber und Upgrades während eines Runs angeboten werden, unterliegt bewusst einem gewissen Zufall. Dadurch entstehen Runs, in denen besonders starke Synergien und funktionierende Builds zustande kommen, ebenso wie Runs, in denen das entwickelte Loadout weniger gut funktioniert.
Entscheidend ist dabei, dass sich die Qualität eines Builds sowohl auf das Spielerlebnis als auch auf die Dauer des Runs auswirkt. Gelingt es, ein besonders starkes und gut synergierendes Loadout aufzubauen, lassen sich Encounter effizienter bewältigen, der Spieler kommt weiter und der Run dauert entsprechend länger. Gerade die Runs, in denen die Systeme besonders gut ineinandergreifen und sich der eigene Build besonders befriedigend anfühlt, nehmen damit automatisch einen größeren Anteil der tatsächlichen Spielzeit ein.
Entsteht dagegen durch ungünstige Belohnungen oder weniger gelungene Build-Entscheidungen ein Loadout, das nur schlecht funktioniert, wird dies zwar ebenfalls spürbar, der betreffende Run endet jedoch typischerweise früher. Der Spieler verbringt dadurch vergleichsweise wenig Zeit in einem Zustand, in dem er bereits merkt, dass sein aktueller Build nicht überzeugend funktioniert, und kann zeitnah mit einem neuen Run und neuen Möglichkeiten beginnen.
Die Run-Struktur erzeugt damit eine günstige zeitliche Gewichtung der durch Zufall und Build-Entscheidungen entstehenden Varianz: Je besser ein Build funktioniert und je befriedigender sich ein Run dadurch anfühlt, desto länger dauert er tendenziell; schwächere Runs verkürzen sich dagegen selbst. Dadurch verbringt der Spieler über die gesamte Spielzeit hinweg überproportional viel Zeit mit den besonders gelungenen und motivierenden Runs.

Auf diese Weise verbindet der Roguelite-Modus unterschiedliche Formen von Fortschritt: Der Spieler kann selbst besser werden, sein Build kann stärker werden und die persistente Metaprogression kann seine zukünftigen Ausgangsbedingungen verbessern. Dadurch können Spieler mit sehr unterschiedlichem Kompetenzniveau über längere Zeit in einem Bereich spielen, der für sie herausfordernd, aber nicht dauerhaft blockierend ist, während gleichzeitig das Gefühl erhalten bleibt, aus eigener Kraft immer weiter in das Spiel vorzudringen.